\subsubsection{Covariance and diversity}
\label{subsec:expression_cov_div}
The covariance term increases when the predictions of $\{f(\cdot, \theta_m)\}_{m=1}^M$ are correlated.
In the worst case where all predictions are identical, covariance equals variance and WA is no longer beneficial.
On the other hand, the lower the covariance, the greater the gain of WA over its members; this is derived by comparing
\Cref{eq:b_v,eq:b_var_cov}, as detailed in \Cref{app:proof_wa_ind}.
It motivates tackling covariance by encouraging members to make different predictions, thus to be functionally diverse.
Diversity is a widely analyzed concept in the ensemble literature \cite{Lakshminarayanan2017}, for which numerous measures have been introduced \cite{kuncheva2003measures,aksela2003comparison,DBLP:journals/corr/abs-1905-00414}.
\reb{In \Cref{sect:dwa}, we aim at decorrelating the learning procedures to increase members' diversity and reduce the covariance term}.

% Deep ensembles \cite{Lakshminarayanan2017} usually seek zero covariance by members' independence, \reb{\eg via explicit adversarial regularization in \cite{rame2021dice}: we implicitly follow the same goal when s in \Cref{sect:dwa}}.
% we follow this objective in \Cref{sect:dwa} by decorrelating the learning procedures.
% For example, \cite{rame2021dice} explicitly sought zero covariance by member's independence in Deep ensembles \cite{Lakshminarayanan2017}.